\chapter{Frobenius algebras and 2D TQFT: full proof}\label{app:frobenius}

This appendix is the landing zone for the parts of Chapter~\ref{ch:frobenius} that were stated honestly but not proved in full in the main text. The guiding principle is simple: every deferred argument in Chapter~\ref{ch:frobenius} should have a named home here, even if the full derivation is postponed until the Stage VI integration pass.

\section{From a Frobenius algebra to a TQFT}\label{appsec:frob-to-tqft}

\emph{To be drafted in Stage VI.}

This section will supply the full converse direction of the classification theorem:
\begin{equation*}
(V,m,\eta,\epsilon)
\quad\Longrightarrow\quad
Z_V : \mathrm{Cob}_2 \to \mathrm{Vect}_{\C}.
\end{equation*}
The main-text version in Proposition~\ref{prop:Frob-to-TQFT} defined $Z_V$ on the generators and explained why the local algebraic relations match the local cobordism relations. What remains here is the global well-definedness argument: given an arbitrary cobordism and a Morse decomposition into caps, cups, cylinders, and pairs of pants, prove that the resulting linear map is independent of the chosen Morse function and chosen sequence of cuts.

\medskip

\noindent Planned contents:
\begin{enumerate}
\item Construct $Z_V$ on generators and disconnected objects.
\item Show invariance under local generator relations.
\item Show that any two Morse decompositions are connected by the standard Cerf-theoretic moves.
\item Conclude that $Z_V$ is a well-defined symmetric monoidal functor.
\end{enumerate}

\section{Independence of handle decomposition and mapping-class-group moves}\label{appsec:mapping-class}

\emph{To be drafted in Stage VI.}

This section will isolate the geometric input behind the converse direction. Chapter~\ref{ch:frobenius} used the slogan that ``the relations already encoded in the Frobenius algebra are enough.'' Here that slogan will be expanded into a precise statement about moves relating decompositions of the same surface.

\medskip

\noindent Planned contents:
\begin{enumerate}
\item The Morse-move/Cerf-move list for surfaces.
\item Why these moves reduce to associativity, commutativity, unitality, and the Frobenius relation.
\item How the resulting invariance packages the mapping class group action on decomposed surfaces.
\item A clean statement of where Kock's proof is being used and where this appendix reproduces the argument directly.
\end{enumerate}

\section{Handle-operator spectrum in the finite-group examples}\label{appsec:handle-spectrum}

\emph{To be drafted in Stage VI.}

Section~\ref{sec:genus-g-partition} derived the general genus formula in terms of the handle operator. The examples in Section~\ref{sec:frobenius-examples} showed how to build Frobenius algebras from $\C[x]/(x^n-1)$ and from $Z(\C[G])$, but did not carry out the full diagonalization of the handle operator in those bases. This section will do that explicitly.

\medskip

\noindent Planned contents:
\begin{enumerate}
\item Compute the handle element and handle operator for $A_n=\C[x]/(x^n-1)$.
\item Compute the handle operator on the class-sum basis of $Z(\C[G])$.
\item Rewrite the genus formula of Theorem~\ref{thm:genus-g-handle} in those examples.
\item Compare the unnormalized Frobenius-theory answer with the conventional finite-gauge-theory normalization.
\end{enumerate}

\medskip

\noindent In particular, this section is the designated home for the statement deferred after Proposition~\ref{prop:idempotent-genus}: the explicit eigenvalues of the handle operator in the finite-group examples used to foreshadow Chapter~\ref{ch:dw}.
